\documentclass[12pt]{article}
\usepackage[utf8]{inputenc}
\usepackage[T1]{fontenc}
\usepackage[margin=1in]{geometry}
\usepackage{setspace}
\usepackage{xcolor}
\usepackage[colorlinks=true, urlcolor=blue!60!black, linkcolor=blue!60!black]{hyperref}
\setlength{\parindent}{0em}
\setlength{\parskip}{0.6em}
\renewcommand{\baselinestretch}{1.0}
\begin{document}
\begin{center}
{\LARGE \textit{West Virginia v.\ EPA}, 597 U.S.\ 697 (2022)\par}
\vspace{0.5em}
{\large Case Brief --- Two-Page Edition\par}
\end{center}

\section{Memory Jogger}
The ``major questions doctrine,'' named for the first time: \S~111(d) did not clearly authorize the Clean Power Plan's shift away from coal (\href{https://supreme.justia.com/cases/federal/us/597/20-1530/}{Justia}).

\section{Detailed Case Facts}
Section 111(d) lets EPA set existing-source limits based on the ``best system of emission reduction.'' EPA had always used it for on-site technology. The 2015 Clean Power Plan instead required ``generation shifting'' from coal to gas and renewables, projecting coal's share of electricity to fall from 38\% to 27\% by 2030.

\section{Procedural History}
The Supreme Court stayed the Plan in 2016; EPA repealed it in 2019 and adopted the narrower ACE Rule. The D.C.\ Circuit vacated both, rejecting the major questions doctrine. The Supreme Court reversed.

\section{Judicial Votes}
Reversed 6--3. Roberts wrote for Thomas, Alito, Gorsuch, Kavanaugh, and Barrett; Gorsuch (with Alito) concurred; Kagan (with Breyer and Sotomayor) dissented.

\section{Holding}
The case was not moot. In extraordinary cases of vast economic and political significance, an agency must point to ``clear congressional authorization''; \S~111(d) supplies none for generation shifting.

\section{Analysis of Opinions}
\textbf{Majority:} EPA claimed a ``transformative expansion'' of power in a statutory ``backwater,'' outside its expertise, implementing a cap-and-trade scheme Congress had ``considered and rejected''; ``system'' is no clear authorization. \textbf{Gorsuch:} the doctrine is a clear-statement rule protecting Article~I. \textbf{Kagan:} the statute's broad text covers the Plan, and the Court's canon is a ``get-out-of-text-free'' card.

\section{Examples: Future Applications}
\textbf{Same-Side:} using a 1974 lab-safety statute to license frontier AI. \textbf{Opposite-Side:} a source-specific, on-site technology standard under \S~111(d). \textbf{Fence-Sitter:} SEC climate-risk disclosure rules.

\section{Critique}
Scholars and Justices debate whether the doctrine is a substantive canon or ordinary contextual textualism, and whether it is a deregulatory ratchet---a charge tested when \textit{Learning Resources v.\ Trump} (2026) struck down a Republican President's tariffs.

\section{Key Quotations}
\begin{quote}
``To convince us otherwise, something more than a merely plausible textual basis for the agency action is necessary. The agency instead must point to `clear congressional authorization' for the power it claims.''
\end{quote}
\noindent --- Chief Justice Roberts, for the Court, 597 U.S.\ at 723.

\end{document}
